\section{Implementation}
\label{sec:implementation}

ECR is open source under the Apache~2.0 licence. It is written in TypeScript, runs on Node.js~22 with no build step, and has one runtime dependency, a schema library. This section describes it at \ecrcommit{}. Sizes are lines that are neither blank nor comments, by the counting rule of the repository.

\textbf{Core and command line.} The schema of a trace and the contract are two files of 85 and 57 lines. The core (337 lines) records a trace through a driver, validates it against the schema and the vocabularies of its adapter, hashes it, resolves parameters, checks a source map, classifies a replay into the four outcomes of Section~\ref{sec:bug}, and guards a repair. It imports no model SDK. The conformance kit is 260 lines. A command line exposes five commands: \code{validate}, \code{compile}, \code{replay}, \code{repair} and \code{conformance}.

\textbf{Reference adapter.} The reference adapter (302 lines) drives a to-do list held in memory and compiles a trace to a test file for the test runner built into Node.js. It implements snapshots and an instrumented copy of its system, so that every law has a subject. Two altered versions of the system serve as fixtures. In the first, an item is no longer found by its title. In the second, a new item starts ticked: for a trace recorded on the original system this is a change of behaviour, and for a trace explored on it this is a system with a defect.

\textbf{Mobile adapter.} The mobile adapter (298 lines) defines nine actions (launch, tap, long press, type, erase, key press, swipe, scroll until visible, hide the keyboard) and two checks (visible, not visible). A step finds its element by accessibility identifier, by visible text or, as a last resort, by a point given in percent of the screen. A press can say where in its element it lands, and a swipe can keep its two ends. Traces compile to Maestro flows~\cite{maestro}: a parameter becomes a variable passed on the command line, and a literal text is escaped because Maestro reads it as a regular expression. In CI, the syntax checker of Maestro validates the compiled flows. Replay runs Maestro. When a flow fails, the target finds the command that failed, in the record that Maestro keeps or in what it prints, and maps it back through the source map: a failed assertion becomes \code{oracle\_failed} with the oracle it checks, and a failed action becomes a divergence on its step, whose type is read from the wording of the message of Maestro. That reading is the weakest part of the target: on a device it has been checked for one cause only, an element that is not found (Sections~\ref{sec:evaluation} and~\ref{sec:case}), and it never reports \code{unexpected\_state}. The adapter has no driver: an external tool explores and writes the trace (Section~\ref{sec:case}).

\textbf{On a device.} A CI job, started by hand, builds the demonstration application of Section~\ref{sec:case}, starts an Android~11 emulator and replays the flows compiled from two traces, with no model. Both traces were written directly from the source, not explored by an agent. In its run of 25~November 2026, the first replay passed. The second stopped on the oracle that its trace records as failed, a counter that reads ``1~tasks'', and ECR reported it as a known bug. It shows that the chain works on a device.

\textbf{Tests, and what is not written.} A suite of 69 tests covers the packages. Two parts are not started: a server through which an agent explores while ECR records each step, and a web adapter with a Playwright target.
